% maestro2tex -- generated master include for ControlAmp.
%% Regenerate with tools/maestro2tex; do not edit by hand.
%%
%% Usage from the TRM (which builds with cwd = latex/TRM/):
%     % maestro2tex -- generated master include for ControlAmp.
%% Regenerate with tools/maestro2tex; do not edit by hand.
%%
%% Usage from the TRM (which builds with cwd = latex/TRM/):
%     % maestro2tex -- generated master include for ControlAmp.
%% Regenerate with tools/maestro2tex; do not edit by hand.
%%
%% Usage from the TRM (which builds with cwd = latex/TRM/):
%     % maestro2tex -- generated master include for ControlAmp.
%% Regenerate with tools/maestro2tex; do not edit by hand.
%%
%% Usage from the TRM (which builds with cwd = latex/TRM/):
%     \input{include/analog/ControlAmp.tex}
%
%% Needs pgfplots, booktabs, siunitx and placeins -- all already in
%% packages-commands.tex. \providecolor needs xcolor (also present).

\providecommand{\MaestroRoot}{include/analog/}

% Categorical palette: slots 1-5 of the validated reference palette,
% in documented order. Dash patterns are a secondary encoding for print.
\providecolor{mtxA}{HTML}{2A78D6}
\providecolor{mtxB}{HTML}{EB6834}
\providecolor{mtxC}{HTML}{1BAF7A}
\providecolor{mtxD}{HTML}{EDA100}
\providecolor{mtxE}{HTML}{E87BA4}

% Keep this block's floats inside this block (see placeins).
\FloatBarrier

\subsection{Potentiostat Control Amplifier}\label{ss:controlamp}

The control amplifier (schematic cell \texttt{anatop\_controlamp}, drawn from
\texttt{AFE\_ECS/folded\_cascode\_pmos}) is amplifier A1 of the rev-2 potentiostat.
It drives the counter electrode directly from its output, takes its feedback from the
reference electrode on the inverting input, and takes its set point from the reference
DAC on the non-inverting input; the cell current is therefore whatever A1 must push
into the cell to hold $V_{\mathrm{RE}}$ at the commanded potential. Unlike the rev-1
front end it drives no current mirror, so both polarities of cell current pass through
it.

The topology is a single-stage folded cascode with a PMOS input pair. A PMOS tail
device feeds the differential pair from \texttt{vdd!}; the pair's drain currents are
folded into two NMOS current sources biased on \texttt{vbn} and pass up through NMOS
cascodes on \texttt{vbnc} into a cascoded PMOS mirror on \texttt{vbpm}/\texttt{vbpc}.
Output and load current are summed at the single high-impedance node where the PMOS and
NMOS cascodes meet, so the amplifier has no internal compensation capacitor --- the load
capacitance at \texttt{vout} is what sets the dominant pole, and every dynamic figure
below is quoted against the \SI{10}{\pico\farad} load the benches apply. A PMOS input
pair puts the input common-mode range against \texttt{vss!} rather than against the
supply, which is what a potentiostat holding a reference electrode near ground needs.
The four cascode gate biases come from an \texttt{anatop\_biasgen\_local} instance
(Section~\ref{ss:biasgen-local}) trimmed to \texttt{BiasAdj}~=~37.

The characterisation below is from schematic-level simulation (no extracted parasitics)
of three views of the \texttt{castalia/anatop\_tb\_controlamp} testbench:
\texttt{TypOpenLoopWBiasTB} (open-loop AC, DC transfer, noise),
\texttt{UnityFeedbackWBiasTB} (loop stability by \texttt{stb} analysis, unity-gain
transfer, noise) and \texttt{SlewOpenLoopWBiasTB} (large-signal step response). Both
noise-driven benches include the local bias generator, so the noise and supply-current
numbers are for the amplifier as biased, not for an idealised cell. This run covers the
single process, temperature and supply point of
Table~\ref{tab:ControlAmp-corners} --- it is a typical-corner design check, not a
corner-signed characterisation.

\input{\MaestroRoot tab_ControlAmp_corners.tex}
\input{\MaestroRoot tab_ctrlamp_dcop.tex}
\input{\MaestroRoot tab_ctrlamp_ac.tex}
\input{\MaestroRoot fig_ctrlamp_ol_gain.tex}
\input{\MaestroRoot fig_ctrlamp_ol_phase.tex}
\input{\MaestroRoot tab_ctrlamp_stability.tex}
\input{\MaestroRoot tab_ctrlamp_swing.tex}
\input{\MaestroRoot tab_ctrlamp_accuracy.tex}
\input{\MaestroRoot fig_ctrlamp_buffer_error.tex}
\input{\MaestroRoot fig_ctrlamp_ol_idd.tex}
\input{\MaestroRoot tab_ctrlamp_largesignal.tex}
\input{\MaestroRoot fig_ctrlamp_slew.tex}
\input{\MaestroRoot fig_ctrlamp_slew_iload.tex}
\input{\MaestroRoot note_ctrlamp_drive.tex}
\input{\MaestroRoot tab_ctrlamp_drivemap.tex}
\input{\MaestroRoot fig_ctrlamp_drivemap.tex}
\input{\MaestroRoot note_ctrlamp_drivemap.tex}
\input{\MaestroRoot tab_ctrlamp_noise.tex}
\input{\MaestroRoot fig_ctrlamp_noise.tex}

\FloatBarrier

%
%% Needs pgfplots, booktabs, siunitx and placeins -- all already in
%% packages-commands.tex. \providecolor needs xcolor (also present).

\providecommand{\MaestroRoot}{include/analog/}

% Categorical palette: slots 1-5 of the validated reference palette,
% in documented order. Dash patterns are a secondary encoding for print.
\providecolor{mtxA}{HTML}{2A78D6}
\providecolor{mtxB}{HTML}{EB6834}
\providecolor{mtxC}{HTML}{1BAF7A}
\providecolor{mtxD}{HTML}{EDA100}
\providecolor{mtxE}{HTML}{E87BA4}

% Keep this block's floats inside this block (see placeins).
\FloatBarrier

\subsection{Potentiostat Control Amplifier}\label{ss:controlamp}

The control amplifier (schematic cell \texttt{anatop\_controlamp}, drawn from
\texttt{AFE\_ECS/folded\_cascode\_pmos}) is amplifier A1 of the rev-2 potentiostat.
It drives the counter electrode directly from its output, takes its feedback from the
reference electrode on the inverting input, and takes its set point from the reference
DAC on the non-inverting input; the cell current is therefore whatever A1 must push
into the cell to hold $V_{\mathrm{RE}}$ at the commanded potential. Unlike the rev-1
front end it drives no current mirror, so both polarities of cell current pass through
it.

The topology is a single-stage folded cascode with a PMOS input pair. A PMOS tail
device feeds the differential pair from \texttt{vdd!}; the pair's drain currents are
folded into two NMOS current sources biased on \texttt{vbn} and pass up through NMOS
cascodes on \texttt{vbnc} into a cascoded PMOS mirror on \texttt{vbpm}/\texttt{vbpc}.
Output and load current are summed at the single high-impedance node where the PMOS and
NMOS cascodes meet, so the amplifier has no internal compensation capacitor --- the load
capacitance at \texttt{vout} is what sets the dominant pole, and every dynamic figure
below is quoted against the \SI{10}{\pico\farad} load the benches apply. A PMOS input
pair puts the input common-mode range against \texttt{vss!} rather than against the
supply, which is what a potentiostat holding a reference electrode near ground needs.
The four cascode gate biases come from an \texttt{anatop\_biasgen\_local} instance
(Section~\ref{ss:biasgen-local}) trimmed to \texttt{BiasAdj}~=~37.

The characterisation below is from schematic-level simulation (no extracted parasitics)
of three views of the \texttt{castalia/anatop\_tb\_controlamp} testbench:
\texttt{TypOpenLoopWBiasTB} (open-loop AC, DC transfer, noise),
\texttt{UnityFeedbackWBiasTB} (loop stability by \texttt{stb} analysis, unity-gain
transfer, noise) and \texttt{SlewOpenLoopWBiasTB} (large-signal step response). Both
noise-driven benches include the local bias generator, so the noise and supply-current
numbers are for the amplifier as biased, not for an idealised cell. This run covers the
single process, temperature and supply point of
Table~\ref{tab:ControlAmp-corners} --- it is a typical-corner design check, not a
corner-signed characterisation.

\input{\MaestroRoot tab_ControlAmp_corners.tex}
\input{\MaestroRoot tab_ctrlamp_dcop.tex}
\input{\MaestroRoot tab_ctrlamp_ac.tex}
\input{\MaestroRoot fig_ctrlamp_ol_gain.tex}
\input{\MaestroRoot fig_ctrlamp_ol_phase.tex}
\input{\MaestroRoot tab_ctrlamp_stability.tex}
\input{\MaestroRoot tab_ctrlamp_swing.tex}
\input{\MaestroRoot tab_ctrlamp_accuracy.tex}
\input{\MaestroRoot fig_ctrlamp_buffer_error.tex}
\input{\MaestroRoot fig_ctrlamp_ol_idd.tex}
\input{\MaestroRoot tab_ctrlamp_largesignal.tex}
\input{\MaestroRoot fig_ctrlamp_slew.tex}
\input{\MaestroRoot fig_ctrlamp_slew_iload.tex}
\input{\MaestroRoot note_ctrlamp_drive.tex}
\input{\MaestroRoot tab_ctrlamp_drivemap.tex}
\input{\MaestroRoot fig_ctrlamp_drivemap.tex}
\input{\MaestroRoot note_ctrlamp_drivemap.tex}
\input{\MaestroRoot tab_ctrlamp_noise.tex}
\input{\MaestroRoot fig_ctrlamp_noise.tex}

\FloatBarrier

%
%% Needs pgfplots, booktabs, siunitx and placeins -- all already in
%% packages-commands.tex. \providecolor needs xcolor (also present).

\providecommand{\MaestroRoot}{include/analog/}

% Categorical palette: slots 1-5 of the validated reference palette,
% in documented order. Dash patterns are a secondary encoding for print.
\providecolor{mtxA}{HTML}{2A78D6}
\providecolor{mtxB}{HTML}{EB6834}
\providecolor{mtxC}{HTML}{1BAF7A}
\providecolor{mtxD}{HTML}{EDA100}
\providecolor{mtxE}{HTML}{E87BA4}

% Keep this block's floats inside this block (see placeins).
\FloatBarrier

\subsection{Potentiostat Control Amplifier}\label{ss:controlamp}

The control amplifier (schematic cell \texttt{anatop\_controlamp}, drawn from
\texttt{AFE\_ECS/folded\_cascode\_pmos}) is amplifier A1 of the rev-2 potentiostat.
It drives the counter electrode directly from its output, takes its feedback from the
reference electrode on the inverting input, and takes its set point from the reference
DAC on the non-inverting input; the cell current is therefore whatever A1 must push
into the cell to hold $V_{\mathrm{RE}}$ at the commanded potential. Unlike the rev-1
front end it drives no current mirror, so both polarities of cell current pass through
it.

The topology is a single-stage folded cascode with a PMOS input pair. A PMOS tail
device feeds the differential pair from \texttt{vdd!}; the pair's drain currents are
folded into two NMOS current sources biased on \texttt{vbn} and pass up through NMOS
cascodes on \texttt{vbnc} into a cascoded PMOS mirror on \texttt{vbpm}/\texttt{vbpc}.
Output and load current are summed at the single high-impedance node where the PMOS and
NMOS cascodes meet, so the amplifier has no internal compensation capacitor --- the load
capacitance at \texttt{vout} is what sets the dominant pole, and every dynamic figure
below is quoted against the \SI{10}{\pico\farad} load the benches apply. A PMOS input
pair puts the input common-mode range against \texttt{vss!} rather than against the
supply, which is what a potentiostat holding a reference electrode near ground needs.
The four cascode gate biases come from an \texttt{anatop\_biasgen\_local} instance
(Section~\ref{ss:biasgen-local}) trimmed to \texttt{BiasAdj}~=~37.

The characterisation below is from schematic-level simulation (no extracted parasitics)
of three views of the \texttt{castalia/anatop\_tb\_controlamp} testbench:
\texttt{TypOpenLoopWBiasTB} (open-loop AC, DC transfer, noise),
\texttt{UnityFeedbackWBiasTB} (loop stability by \texttt{stb} analysis, unity-gain
transfer, noise) and \texttt{SlewOpenLoopWBiasTB} (large-signal step response). Both
noise-driven benches include the local bias generator, so the noise and supply-current
numbers are for the amplifier as biased, not for an idealised cell. This run covers the
single process, temperature and supply point of
Table~\ref{tab:ControlAmp-corners} --- it is a typical-corner design check, not a
corner-signed characterisation.

\input{\MaestroRoot tab_ControlAmp_corners.tex}
\input{\MaestroRoot tab_ctrlamp_dcop.tex}
\input{\MaestroRoot tab_ctrlamp_ac.tex}
\input{\MaestroRoot fig_ctrlamp_ol_gain.tex}
\input{\MaestroRoot fig_ctrlamp_ol_phase.tex}
\input{\MaestroRoot tab_ctrlamp_stability.tex}
\input{\MaestroRoot tab_ctrlamp_swing.tex}
\input{\MaestroRoot tab_ctrlamp_accuracy.tex}
\input{\MaestroRoot fig_ctrlamp_buffer_error.tex}
\input{\MaestroRoot fig_ctrlamp_ol_idd.tex}
\input{\MaestroRoot tab_ctrlamp_largesignal.tex}
\input{\MaestroRoot fig_ctrlamp_slew.tex}
\input{\MaestroRoot fig_ctrlamp_slew_iload.tex}
\input{\MaestroRoot note_ctrlamp_drive.tex}
\input{\MaestroRoot tab_ctrlamp_drivemap.tex}
\input{\MaestroRoot fig_ctrlamp_drivemap.tex}
\input{\MaestroRoot note_ctrlamp_drivemap.tex}
\input{\MaestroRoot tab_ctrlamp_noise.tex}
\input{\MaestroRoot fig_ctrlamp_noise.tex}

\FloatBarrier

%
%% Needs pgfplots, booktabs, siunitx and placeins -- all already in
%% packages-commands.tex. \providecolor needs xcolor (also present).

\providecommand{\MaestroRoot}{include/analog/}

% Categorical palette: slots 1-5 of the validated reference palette,
% in documented order. Dash patterns are a secondary encoding for print.
\providecolor{mtxA}{HTML}{2A78D6}
\providecolor{mtxB}{HTML}{EB6834}
\providecolor{mtxC}{HTML}{1BAF7A}
\providecolor{mtxD}{HTML}{EDA100}
\providecolor{mtxE}{HTML}{E87BA4}

% Keep this block's floats inside this block (see placeins).
\FloatBarrier

\subsection{Potentiostat Control Amplifier}\label{ss:controlamp}

The control amplifier (schematic cell \texttt{anatop\_controlamp}, drawn from
\texttt{AFE\_ECS/folded\_cascode\_pmos}) is amplifier A1 of the rev-2 potentiostat.
It drives the counter electrode directly from its output, takes its feedback from the
reference electrode on the inverting input, and takes its set point from the reference
DAC on the non-inverting input; the cell current is therefore whatever A1 must push
into the cell to hold $V_{\mathrm{RE}}$ at the commanded potential. Unlike the rev-1
front end it drives no current mirror, so both polarities of cell current pass through
it.

The topology is a single-stage folded cascode with a PMOS input pair. A PMOS tail
device feeds the differential pair from \texttt{vdd!}; the pair's drain currents are
folded into two NMOS current sources biased on \texttt{vbn} and pass up through NMOS
cascodes on \texttt{vbnc} into a cascoded PMOS mirror on \texttt{vbpm}/\texttt{vbpc}.
Output and load current are summed at the single high-impedance node where the PMOS and
NMOS cascodes meet, so the amplifier has no internal compensation capacitor --- the load
capacitance at \texttt{vout} is what sets the dominant pole, and every dynamic figure
below is quoted against the \SI{10}{\pico\farad} load the benches apply. A PMOS input
pair puts the input common-mode range against \texttt{vss!} rather than against the
supply, which is what a potentiostat holding a reference electrode near ground needs.
The four cascode gate biases come from an \texttt{anatop\_biasgen\_local} instance
(Section~\ref{ss:biasgen-local}) trimmed to \texttt{BiasAdj}~=~37.

The characterisation below is from schematic-level simulation (no extracted parasitics)
of three views of the \texttt{castalia/anatop\_tb\_controlamp} testbench:
\texttt{TypOpenLoopWBiasTB} (open-loop AC, DC transfer, noise),
\texttt{UnityFeedbackWBiasTB} (loop stability by \texttt{stb} analysis, unity-gain
transfer, noise) and \texttt{SlewOpenLoopWBiasTB} (large-signal step response). Both
noise-driven benches include the local bias generator, so the noise and supply-current
numbers are for the amplifier as biased, not for an idealised cell. This run covers the
single process, temperature and supply point of
Table~\ref{tab:ControlAmp-corners} --- it is a typical-corner design check, not a
corner-signed characterisation.

\input{\MaestroRoot tab_ControlAmp_corners.tex}
\input{\MaestroRoot tab_ctrlamp_dcop.tex}
\input{\MaestroRoot tab_ctrlamp_ac.tex}
\input{\MaestroRoot fig_ctrlamp_ol_gain.tex}
\input{\MaestroRoot fig_ctrlamp_ol_phase.tex}
\input{\MaestroRoot tab_ctrlamp_stability.tex}
\input{\MaestroRoot tab_ctrlamp_swing.tex}
\input{\MaestroRoot tab_ctrlamp_accuracy.tex}
\input{\MaestroRoot fig_ctrlamp_buffer_error.tex}
\input{\MaestroRoot fig_ctrlamp_ol_idd.tex}
\input{\MaestroRoot tab_ctrlamp_largesignal.tex}
\input{\MaestroRoot fig_ctrlamp_slew.tex}
\input{\MaestroRoot fig_ctrlamp_slew_iload.tex}
\input{\MaestroRoot note_ctrlamp_drive.tex}
\input{\MaestroRoot tab_ctrlamp_drivemap.tex}
\input{\MaestroRoot fig_ctrlamp_drivemap.tex}
\input{\MaestroRoot note_ctrlamp_drivemap.tex}
\input{\MaestroRoot tab_ctrlamp_noise.tex}
\input{\MaestroRoot fig_ctrlamp_noise.tex}

\FloatBarrier
